\section{Proof of Cubic-Quintic Scheme}
\label{sec:certification}

Write the correlation function of the cubic--quintic limiting scheme as
\begin{equation}
  H(t)=b_1t+\sum_{\substack{m\ge 3\\ m\ \mathrm{odd}}} b_m t^m,
  \label{eq:certification-H-series}
\end{equation}
and let
\begin{equation}
  \Delta_H:=\sum_{\substack{m\ge 3\\ m\ \mathrm{odd}}}|b_m|
  \label{eq:nonlinear-coefficient-mass}
\end{equation}
be the total mass of its nonlinear coefficients.  Let
\[
  G(z)=H^{-1}(z)=\sum_{n\ge1}a_nz^n
\]
denote the local holomorphic inverse of $H$ at the origin, and define its
majorant by
\begin{equation}
  \mathcal M(r):=\sum_{n\ge1}|a_n|r^n.
  \label{eq:inverse-majorant-certification}
\end{equation}
The local inverse exists because $H'(0)=b_1>0$, by the holomorphic inverse
function theorem.

For the linear map $t\mapsto b_1t$, the condition $\mathcal M(\gamma)<1$ is
simply $\gamma<b_1$.  The next proposition shows that the nonlinear part costs
at most $\Delta_H$ of this margin.

\begin{proposition}[Inverse certificate from near-linearity]
\label{prop:inverse-certificate-near-linearity}
Assume $b_1>0$ and $\Delta_H<\infty$.  If
\begin{equation}
  \gamma+\Delta_H<b_1,
  \label{eq:near-linearity-condition}
\end{equation}
then there exists $R>\gamma$ such that the local inverse $G$ extends
holomorphically to $R\mathbb D$, maps $R\mathbb D$ into $\mathbb D$, and
satisfies
\begin{equation}
  \mathcal M(\gamma)
  \le \frac{\gamma+\Delta_H}{b_1}
  <1.
  \label{eq:near-linearity-conclusion}
\end{equation}
\end{proposition}

\begin{proof}
Set
\[
  E(t):=\sum_{m\ge3}b_mt^m,
  \qquad
  E_{\mathrm{maj}}(u):=\sum_{m\ge3}|b_m|u^m.
\]
There is a unique formal power series
\[
  V(z)=\sum_{n\ge1}v_nz^n
\]
satisfying
\begin{equation}
  b_1V(z)=z+E_{\mathrm{maj}}(V(z)).
  \label{eq:majorizing-inverse-equation}
\end{equation}
Indeed, comparison of the coefficient of $z$ gives $v_1=b_1^{-1}$.  For
$n\ge2$, the coefficient of $z^n$ in $E_{\mathrm{maj}}(V)$ depends only on
$v_1,\ldots,v_{n-1}$, so \eqref{eq:majorizing-inverse-equation} determines
$v_n$ recursively.  The same recursion shows that $v_n\ge0$ for every $n$.

Comparing coefficients in
\[
  b_1G(z)=z-E(G(z))
\]
and in \eqref{eq:majorizing-inverse-equation}, induction on $n$ gives
\begin{equation}
  |a_n|\le v_n
  \qquad(n\ge1).
  \label{eq:inverse-coefficient-domination}
\end{equation}
Indeed, at degree $n$, every nonlinear composition involves only previously
determined coefficients, and replacing each $b_m$ by $|b_m|$ removes all
possible cancellation.

Choose $R>\gamma$ so close to $\gamma$ that
\[
  q:=\frac{R+\Delta_H}{b_1}<1.
\]
Starting from $V_0(z)=0$, define
\begin{equation}
  V_{k+1}(z):=\frac{z+E_{\mathrm{maj}}(V_k(z))}{b_1}.
  \label{eq:majorizing-inverse-iteration}
\end{equation}
The series $V_k$ have nonnegative coefficients and increase coefficientwise,
meaning that each Taylor coefficient is nondecreasing in $k$.
Moreover, if $V_k(R)\le q$, then $V_k(R)<1$ and
\[
  V_{k+1}(R)
  \le \frac{R+E_{\mathrm{maj}}(1)}{b_1}
  =q.
\]
Thus $V_k(R)\le q$ for every $k$.  Their coefficientwise limit satisfies
\eqref{eq:majorizing-inverse-equation}, and hence equals $V$ by uniqueness.
The monotone convergence theorem now gives
\[
  V(R)=\sum_{n\ge1}v_nR^n\le q<1.
\]
Together with \eqref{eq:inverse-coefficient-domination}, this implies
\[
  \sum_{n\ge1}|a_n|R^n\le V(R)<1.
\]
Consequently, the Taylor series of $G$ converges on $R\mathbb D$ and takes
values in $\mathbb D$.  Since the identity $H(G(z))=z$ holds near the origin,
it holds throughout $R\mathbb D$ by the identity theorem.  Finally,
\[
  \mathcal M(\gamma)
  \le V(\gamma)
  =\frac{\gamma+E_{\mathrm{maj}}(V(\gamma))}{b_1}
  \le\frac{\gamma+\Delta_H}{b_1},
\]
because $0\le V(\gamma)\le V(R)<1$.
\end{proof}

It remains to certify $\Delta_H$.  We compute the coefficients through a
finite degree $N$ and bound the remaining tail indirectly.  The point is that
the Euler operator
\[
  D:=t\frac{d}{dt}
\]
multiplies the coefficient of $t^m$ by $m$, so $D^3$ strongly amplifies
high-degree coefficients.  The $L^2(\mathbb T)$ norm below is the
root-mean-square size of the resulting function around the unit circle
$\mathbb T:=\{z\in\mathbb C:|z|=1\}$.

\begin{lemma}[Third-order coefficient-tail bound]
\label{lem:third-order-tail}
Suppose that $D^3H$ has radial boundary values in $L^2(\mathbb T)$, with
\[
  \|D^3H\|_{L^2(\mathbb T)}^2
  :=\frac{1}{2\pi}\int_0^{2\pi}
  |D^3H(e^{i\theta})|^2\,d\theta.
\]
Then, for every odd $N$,
\begin{equation}
  \sum_{\substack{m>N\\m\ \mathrm{odd}}}|b_m|
  \le
  \|D^3H\|_{L^2(\mathbb T)}
  \left(\sum_{\substack{m>N\\m\ \mathrm{odd}}}m^{-6}\right)^{1/2}
  \le
  \frac{\|D^3H\|_{L^2(\mathbb T)}}{\sqrt{10N^5}}.
  \label{eq:third-order-tail}
\end{equation}
\end{lemma}

\begin{proof}
Since $D^3(t^m)=m^3t^m$, Parseval's identity gives
\begin{equation}
  \|D^3H\|_{L^2(\mathbb T)}^2
  =\sum_{\substack{m\ge1\\m\ \mathrm{odd}}}m^6|b_m|^2.
  \label{eq:parseval-third-order}
\end{equation}
The Cauchy--Schwarz inequality therefore gives
\[
  \sum_{\substack{m>N\\m\ \mathrm{odd}}}|b_m|
  =\sum_{\substack{m>N\\m\ \mathrm{odd}}}
  m^{-3}\bigl(m^3|b_m|\bigr)
  \le
  \left(\sum_{\substack{m>N\\m\ \mathrm{odd}}}m^{-6}\right)^{1/2}
  \|D^3H\|_{L^2(\mathbb T)}.
\]
Finally, since $x\mapsto x^{-6}$ is decreasing and consecutive odd integers
are separated by two, the monotone integral estimate gives
\[
  \sum_{\substack{m>N\\m\ \mathrm{odd}}}m^{-6}
  \le\frac12\int_N^\infty x^{-6}\,dx
  =\frac{1}{10N^5}.
\]
\end{proof}

\begin{proposition}[Certified numerical input]
\label{prop:certified-numerical-input}
For the cubic--quintic correlation function, interval arithmetic certifies
\begin{equation}
  b_1\ge0.881573822049,
  \qquad
  \sum_{\substack{3\le m\le251\\m\ \mathrm{odd}}}|b_m|
  \le1.1328860\cdot10^{-5},
  \qquad
  \|D^3H\|_{L^2(\mathbb T)}<14.443.
  \label{eq:certified-numerical-input}
\end{equation}
\end{proposition}

\begin{proof}
See Appendix~\ref{app:certification-details}.
\end{proof}

Lemma~\ref{lem:third-order-tail} and
Proposition~\ref{prop:certified-numerical-input} give
\begin{equation}
  \Delta_H
  \le
  1.1328860\cdot10^{-5}
  +\frac{14.443}{\sqrt{10\cdot251^5}}
  <1.5904724\cdot10^{-5}.
  \label{eq:certified-nonlinear-mass}
\end{equation}

\begin{theorem}[Cubic--quintic upper bound]
\label{thm:cubic-quintic-upper-bound-certified}
For
\[
  \gamma=0.881545409,
\]
the cubic--quintic limiting scheme satisfies $\mathcal M(\gamma)<1$.
Consequently,
\begin{equation}
  K_G\le\frac{\pi}{2\gamma}
  \le1.7818666069360661.
  \label{eq:certified-KG-bound}
\end{equation}
\end{theorem}

\begin{proof}
By \eqref{eq:certified-nonlinear-mass},
\[
  \gamma+\Delta_H
  <0.881545409+1.5904724\cdot10^{-5}
   =0.881561313724
   <0.881573822049
   \le b_1
\]
Proposition~\ref{prop:inverse-certificate-near-linearity} therefore gives
$\mathcal M(\gamma)<1$, together with the analytic room required by the
finite-dimensional transfer theorem.  That theorem passes the strict
certificate to all sufficiently large ordinary finite-dimensional
approximations.  The classical Krivine preprocessing theorem then gives
\eqref{eq:certified-KG-bound}.
\end{proof}


\paragraph{Reproducibility.}
  The scripts, certificates, and Arb outputs
  are available at \url{https://github.com/trishullab/grothendieck-bounds}. 